\scanpage{8}
Dans le cas général
\[
I^n=0
\]
on procède par récurrence sur $n$, en notant que c'est trivial pour $n=\unsure{1|0}$, déjà fait pour $n=2$, on peut supposer $n>2$ et la chose \deleted{faite} pour $n'<n$. \deleted{\unsure{Cela|Il}} Il suffit alors d'utiliser l'hyp de récurrence et le

\uline{Lemme} \quad Si $I,J$ sont deux idéaux de $A$ tels que $\lambda\mapsto\lambda 1_B$ est une bijection de $I$ ($J$) sur un idéal de $B$, alors $IJ$ a la $=$ propriété, donc aussi les puissances \inserted{$I^n$} de $I$, \unsure{avec} $n\geq1$. \deleted{Si de plus}

$+$ \quad (lemme trivial)

\uline{Corollaire} 1 \quad Soient $A,I,B\deleted{=A'}$ comme dessus [$A\subset B$, $I\subset A$ un idéal \uline{nilpotent} de $A$ tel que $IB=\unsure{B|I}$] \deleted{et tel que} supposons $B_0/A_0$ \deleted{fid. plat} \inserted{morphisme de descente effective pour la cat. \deleted{\unsure{fib}}fibrée des modules plats, \uline{p.ex.} $B_0/A_0$ fid plat}. Alors $\operatorname{Spec}B\longrightarrow\operatorname{Spec}A$ est un morphisme de descente \uline{effective} pour la catégorie fibrée des modules \uline{plats}.
% archive sheet number: 7
